
% !TEX root = ../main.tex
% Resumen del Primer Avance - Proyecto Campus Zero

\section{Antecedentes e Integración del Primer Avance}
\label{sec:avance1}

\subsection{Problema y solución}

En las salas de estudio de la UTFSM, Casa Central, las luces y extractores de aire pueden permanecer encendidos después de que los usuarios se retiran, generando un consumo eléctrico innecesario. Considerando una estimación inicial de tres horas diarias sin ocupantes, una potencia de 0,26~kW por sala y 22 días hábiles al mes, el desperdicio alcanzaría 85,8~kWh mensuales en cinco salas. Esta cifra deberá comprobarse mediante mediciones.

Para enfrentar el problema, \textit{Campus Zero} propone instalar un sistema automatizado en al menos cinco salas, utilizando sensores de presencia infrarrojos (PIR), relés inteligentes y microcontroladores. El sistema apagará las luces y extractores tras diez minutos continuos sin detectar personas, sin intervenir los computadores. Además, se comparará el consumo antes y después de la instalación para evaluar el ahorro obtenido.

\subsection{Beneficiarios y objetivos}

La beneficiaria directa es la UTFSM, particularmente la administración de Casa Central, por la posibilidad de reducir sus gastos eléctricos y el impacto ambiental asociado. Los beneficiarios indirectos son estudiantes, docentes y otros usuarios de las salas, con una capacidad estimada de hasta 40 personas simultáneas.

\noindent\textbf{Objetivo general:} Reducir el consumo eléctrico innecesario de iluminación y extractores de aire mediante un sistema piloto de control automatizado en al menos cinco salas durante 2026, estableciendo una solución replicable para la universidad.

\noindent\textbf{Objetivos específicos:}
\begin{enumerate}
    \item \textbf{OE1:} Diseñar el sistema y adquirir los componentes para al menos cinco salas en un máximo de cinco semanas desde la adjudicación de los fondos, sin superar \$1.250.000.

    \item \textbf{OE2:} Instalar y calibrar el sistema en las cinco salas en un máximo de cuatro semanas desde la recepción de los materiales, configurando el apagado tras diez minutos continuos sin presencia.

    \item \textbf{OE3:} Medir y comparar el consumo eléctrico antes y después de la instalación durante cuatro semanas de prueba, buscando una reducción de al menos un 50\% del consumo durante los períodos de desocupación.

    \item \textbf{OE4:} Entregar el informe de resultados, rendición de cuentas, inventario de equipos y guía de replicación dentro de las cuatro semanas posteriores a la finalización de las mediciones.
\end{enumerate}

\subsection{Financiamiento y continuidad}

El proyecto pertenece a la línea de \textbf{Sustentabilidad y Medio Ambiente}, con un presupuesto máximo de \$1.250.000 para implementar el piloto. En este Segundo Avance se presenta la Carta Gantt, el presupuesto detallado y la red de colaboradores necesarios para planificar su ejecución y evaluar su viabilidad.

